\cotefile{156-9}
\scanpage{10}
\mspage{8}
% source: top, two-line display
\[
\left\{\begin{aligned}
I_{<a}\overset{\unsure{\circ|0}}{\ll}I_{<b}&\Longleftrightarrow a\leq b\\
X\overset{\circ}{\ll}I_{<a}&\deleted{\unsure{\&}}\Longleftrightarrow X\deleted{\unsure{<}}\leq a
\end{aligned}\right.
\qquad \deleted{\unsure{\text{et}}}\quad\text{(cf déf\ldots)}
\]
% source: upper-left margin
% margin note written sideways; slanted
\margin{On \uuline{\unsure{a plus|de plus}} $X\overset{\circ}{\ll}Y\Longleftrightarrow |X|\subset|Y|$ et $X\unsure{\overline{\triangleleft}|\overline{\triangle}}Y$ pour $X,Y\in\unsure{\widehat M|\widehat{\mathscr M}}$ et \unsure{remarquer} \deleted{\unsure{(*)}}\unsure{?} C'est vrai \unsure{seult|seulement} si $\unsure{L|\mathscr L}$ n'a ni plus petit ni plus grand élément.}
% source: top, continuation below display
(où $X\in M$), et \unsure{aussi} pour \deleted{$\unsure{I_{<a}}$} les relations \deleted{\unsure{des semi}} impliquant $I_{>a}$ \unsure{intervalles} \unsure{infinis|définis} symétriquement.

On note.

% source: upper-middle, a)
a) On n'a jamais une relation
\[
I_{<a}\overset{\circ}{\ll}X\qquad\text{pour }X\in M
\]
même si $L$ a un plus petit él. $\unsure{a_0|a_o}$.

% source: middle, example
Par exemple, \deleted{\unsure{que nous}} on n'a \uline{pas} $I_{<a}\overset{\circ}{\ll}I_{\omega_0,a}$, \deleted{car (\unsure{la même heuristique},} \unsure{donc pas} $I_{\omega_0,a}\overset{\circ}{\ll}I_{<a}$. En fait, $I_{\omega_0,a}$ doit être vu (en tant que définissant une \unsure{multiplicité}\ldots) comme une \uline{subdivision} effective de $I_{<a}$.

(Si on avait les deux relations à la fois, on \deleted{aurait}\inserted{devrait avoir} $I_{<a}=I_{\omega_0,a}$, ce qui n'est pas conforme aux intuitions \unsure{présidant} aux \unsure{introductions} de $I_{<a}$.
% source: small added paragraph below preceding lines
\inserted{\unsure{Bien sûr}, \unsure{on aura|on n'aura} $\deleted{\unsure{\not\ll}}I_{<a}\ll I_{\omega_0,a}$, mais la relation inverse \deleted{\unsure{seule}} $I_{\omega_0,a}\ll I_{<a}$, est vraie malgré que $|I_{\omega_0,a}|=|I_{<a}|$.}

% source: lower-middle, b)
b) On n'a jamais une relation
\[
I_{<a}\overset{\circ}{\ll}I_{>b},\quad\text{ou l'inverse.}
\]
% source: lower-left margin
% slanted
\margin{\unsure{Ne plus|Ne pas}\unsure{?} sur $\widehat M$}
% source: bottom, cancelled definition
Relation d'ordre $\leq$ sur $\widehat M$. On \unsure{pose}, \inserted{\unsure{encore},} \deleted{si $\unsure{X,Y\in}\ X,Y\in\widehat M$} \unsure{\text{Ord}} $\leq:\widehat M$ \unsure{par}
\[
\deleted{X\leq Y\Longleftrightarrow\forall x\in\delta X,\ \forall y\in\delta Y,\quad x\leq y}
\]
\deleted{$\unsure{I_{<}}$}
\deleted{(on écrit aussi $\delta X\leq\delta Y$ relation qui a un sens pour deux parties $A,B$ quelc de $M$: $A\leq B$)}
% archive sheet number: 9
